% !TeX root = ./../main.tex

\section{Charakteryzacja i modelowanie szumu urządzeń}\label{sec:lr:noise}
\agentinfo{
Stan: treść napisana w zadaniu PHD-TS-34 (2026-11-27), poprawiona po recenzji (PHD-TS-90), uzupełniona o akapit par:lr:noise:generators (PHD-TS-94) i skrócona w zadaniu PHD-TS-96 (lista cięć: \agentref{eksperymenty/recenzje/rozdzial-3-skracanie.md}); czeka na przegląd autora.
Notatka źródłowa (tabele prac, zestawienie założeń PDU z literaturą, zakres wyszukiwania, pozycje rezerwowe): \agentref{eksperymenty/przeglad-literatury/sekcja-3.5-szum.md}; teksty źródeł: \agentref{eksperymenty/przeglad-literatury/_robocze/PHD-TS-21/txt}.
Uzgodnienia: definicji metod charakteryzacji nie powtarzać (podrozdział 2.2); wnioski dla PDU przedstawiać jako pośrednie, bo prace mierzą parametry składowe, a nie systematyczny błąd wyniku obwodu; dowody przeciw stałości błędu do następnej kalibracji podawać bez łagodzenia.
Ilorazy i zmiany procentowe dla pracy yeter-aydeniz_measuring_2023 są obliczeniami własnymi z tabel C2 i C3 wersji arXiv 2201.02899.
Sekcja jest tłem rozdziałów 4 (założenia Z1--Z5 i Z8; rozdział 4 odwołuje się do tej sekcji, do tabeli tab:lr:noise:time i do akapitu par:lr:noise:generators), 7 i 8 (rozbieżność symulatora i sprzętu).
Akapit par:lr:noise:generators: decyzją autora z 2026-11-27 dawny rozdział 9 o strukturze szumu nie jest osobnym rozdziałem, a omówienie pracy raubitzek_quantum_noise jest tutaj (plik archiwalny: \agentref{chapters/noise_structure.tex}; recenzja: \agentref{eksperymenty/recenzje/rozdzial-9.md}); to praca powiązana z udziałem autora rozprawy, a nie wkład rozprawy (udział autora: sekcja sec:introduction:publications).
Źródła akapitu: wersja opublikowana (APS Open Science 1, 000109, 2026; tekst w \agentref{eksperymenty/_build/PHD-TS-38/work/pub/published.txt}) i wersja zgłoszeniowa \agentref{artykuły/quantum_noise_elsevier_submission/main.tex}; pierwszeństwo ma wersja opublikowana; daty przebiegów pochodzą z wykazu plików repozytorium danych pracy, nie z pracy.
Ustalenie B1 recenzji (konwencja zapisu macierzy z urządzenia; niepotwierdzone przez autora i współautorów): akapit nie podaje zliczeń elementów zgodnych, wartości wskaźnika zgodności ani obliczeń własnych na elementach spoza przekątnej bloku unitalnego, nie rozstrzyga, która część macierzy jest wierszem afinicznym, a która wektorem przesunięcia, i referuje wyniki z przypisaniem autorom; w tekście drukowanym nie pisać, że opublikowana praca zawiera błąd.
Akapit nie proponuje nowych eksperymentów i nie nazywa błędu stabilną cechą urządzenia.
W zadaniu PHD-TS-96 akapit skrócono mniej więcej o połowę: wypadły szczegóły modelu (po jednym zaburzeniu zespolonym 2x2 na każdą z trzech bramek natywnych, ograniczenie modułu parametrów, model w przód z 12 obwodami tomografii, brak narzuconego zachowania śladu i całkowitej dodatniości, 1000 symulowanych rekonstrukcji), trzy kroki schematu wykorzystania opisu w mitygacji, porównanie zespołów losowych zaburzeń, zestawienie z taksonomią blume-kohout_taxonomy_2022 i zdanie o tym, że praca nie podaje numeru kubitu ani liczby powtórzeń (jest w pytaniu 4 niżej i w dodatku A); zostały dane, rozbieżność co do okresu zbierania danych, wynik na przekątnej, ograniczenia i warunek stabilności schematu.
Z tabeli tab:lr:noise:time usunięto w zadaniu PHD-TS-96 wiersze prac klimov_fluctuations_2018 i carroll_dynamics_2022 (obie zostają cytowane w tekście).
}
\agenttask{
Do decyzji autora (PHD-TS-34).
(1) Akapit o znaczeniu dla założeń PDU zawiera oceny własne: literatura nie potwierdza założenia Z4 w postaci mocnej, linia trendu może nie wykryć zmian epizodycznych, reguła jednej kalibracji PDU na dobę zakłada kalibrację urządzenia jako jedno znane zdarzenie. Czy te oceny mają zostać w przeglądzie, czy przejść do dyskusji ograniczeń w rozdziale 4 (sekcja 4.8), a tutaj zostać tylko zestawienie prac.
(2) Czy w rozdziale 4 opatrzyć zastrzeżeniem zdanie o jednej kalibracji urządzenia na dobę (akapit Warunek ponownej kalibracji): według etxezarreta_multiqubit_2023 urządzenia santiago i bogota kalibrowano kilka razy na dobę, a według wilson_just-in-time_2020 bywały okresy 28--49 godzin bez kalibracji; obserwacja z pracy wilson_just-in-time_2020 pochodzi z okresu wcześniejszego niż eksperymenty PDU, a praca etxezarreta_multiqubit_2023 nie podaje roku pomiarów (dotyczy tych samych procesorów), więc nie wolno pisać, że to inne okresy. Zastrzeżenie do reguły dopisano w podrozdziale 4.8 (zadanie PHD-TS-57, 2026-11-27); do decyzji autora zostaje złagodzenie samych zdań w podrozdziale 4.5.
(3) Hipoteza do sprawdzenia w danych autora, niewpisana do tekstu: większy rozrzut błędu na urządzeniu Santiago niż na urządzeniu Manila może mieć związek z częstszymi kalibracjami tego urządzenia.
Wartości liczbowej dla urządzenia Santiago punkt nie przesądza, bo źródła podają ją w odwrotnym znaczeniu: według artykułu ICCS ponad 56 procent wartości błędu odbiegało od średniej o więcej niż 20 procent, a według pracy magisterskiej 56 procent wartości mieściło się w tym paśmie; dla urządzenia Manila oba źródła podają 83 procent wartości w paśmie (rozbieżność jest zgłoszona w recenzji rozdziału 4, ustalenie B1, i rozstrzyga ją autor albo dane pierwotne).
(4) Czy dodać pozycje rezerwowe z notatki (Schlor 2019, Dahlhauser i Humble 2021, Chen 2023), dziś niecytowane.
}

Każda korekta wyznaczona z~wcześniejszego pomiaru, od macierzy odczytu po funkcję PDU i~modele uczone (podrozdział~\ref{sec:lr:ml}), zakłada, że zmierzona charakterystyka opisuje urządzenie także w~chwili obliczenia; istotne jest więc, czego metody charakteryzacji (podrozdział~\ref{subsec:tb:noise:characterisation}) nie obejmują i~jak szybko szum się zmienia.

Jak pokazano w~pracy~\cite{Merkel2013}, tomografia procesu przypisuje operacji także błędy przygotowania stanu i~pomiaru; tomografia zbioru bramek jest od tej wady wolna~\cite{BlumeKohout2017,bib:quantum_gate_tomography}, zakłada jednak szum stały w~czasie zbierania danych, które może trwać kilka godzin~\cite{hashim_practical_2025}.
% ŹRÓDŁO: hashim_practical_2025, część o ograniczeniach eksperymentalnych (dwukubitowa tomografia zbioru bramek na platformach atomowych), arXiv:2408.12064v2; numeracji sekcji w wersji PRX Quantum nie sprawdzano.
Test z~losowymi sekwencjami bramek skaluje się lepiej, lecz zwraca jedną liczbę, która nie rozróżnia błędów koherentnych i~niekoherentnych~\cite{Magesan2011,ProctorWhat2017,Wallman2015,Hashim2023}, a~modele uczone w~większej skali opisują szum uśredniony do kanału Pauliego~\cite{harper_efficient_2020,van_den_berg_probabilistic_2023}.
% ŹRÓDŁO: Hashim2023, streszczenie wersji opublikowanej (Crossref, DOI 10.1038/s41534-023-00764-y); harper_efficient_2020, streszczenie, arXiv:1907.13022; van_den_berg_probabilistic_2023, wstęp, arXiv:2201.09866.
Poza zasięgiem modeli lokalnych i~markowowskich pozostają przesłuch~\cite{sarovar_detecting_2020,mundada2019crosstalk}, wyciek~\cite{Wood2015} i~pamięć otoczenia~\cite{white_demonstration_2020}, czego skutkiem jest rozbieżność między przewidywaniem z~danych kalibracyjnych a~wynikiem obwodu.
Na dwunastu procesorach programy o~regularnej strukturze zawodziły przy rozmiarze nawet o~rząd wielkości mniejszym niż programy losowe~\cite{proctor_measuring_2022}.
Niewierność bramek CNOT wyznaczona metodą \emph{cycle benchmarking}, czyli dla całego cyklu operacji na rejestrze, była od $1{,}3$ do $4{,}0$ razy większa od wartości kalibracyjnych producenta (obliczenie własne z~tabel pracy~\cite{yeter-aydeniz_measuring_2023}; metody nie mierzą tej samej wielkości).
% ŹRÓDŁO: proctor_measuring_2022, streszczenie, arXiv:2008.11294; yeter-aydeniz_measuring_2023, wstęp (cycle benchmarking: błąd cyklu operacji wykonywanych na rejestrze, ,,when a clock cycle of operations is applied to a specific quantum register'') oraz tabele C2 i C3, arXiv:2201.02899 (ar5iv): ilorazy wartości z cycle benchmarking i z danych producenta dla 18 par liczb od 1,34 do 4,04 (obliczenie własne, PHD-TS-34).

\paragraph{Zmienność szumu w~czasie.}
Pomiary zestawia według skali czasu tabela~\ref{tab:lr:noise:time}.
W~kubitach nadprzewodzących czas relaksacji $\Tone$ zmienia się wskutek niestabilności defektów materiałowych sprzężonych z~kubitem, przez skokowe przełączenia, a~nie powolny trend~\cite{klimov_fluctuations_2018,burnett_decoherence_2019,carroll_dynamics_2022}.
% ŹRÓDŁO: burnett_decoherence_2019, streszczenie i część ,,T1 decay-profiles'', arXiv:1901.04417 (ar5iv); klimov_fluctuations_2018, tekst przy rys. 2, wersja arXiv (PDF).
Dla rozprawy najważniejsze są pomiary wewnątrz okna kalibracji na pięciu pięciokubitowych procesorach IBM używanych także w~rozdziale~\ref{chapter:pdu}: współczynnik zmienności $\Tone$ poszczególnych kubitów wynosił od około $5\%$ do $47\%$, a~fluktuacje nie były istotnie skorelowane między kubitami~\cite{etxezarreta_multiqubit_2023}.
% ŹRÓDŁO: etxezarreta_multiqubit_2023, tabele II i III, sekcja III-A, dodatek C, arXiv:2207.06838 (ar5iv).
Zmieniają się także błędy bramek, odczyt, dokładność obwodów i, w~skali miesięcy, rozkłady parametrów szumu (tabela~\ref{tab:lr:noise:time}), a~dryf obserwowano również w~pułapce jonowej~\cite{proctor_detecting_2020}.

\begin{table}[tb]
\centering
\caption[Pomiary zmienności szumu w~czasie]{Pomiary zmienności szumu w~czasie według skali; poza dokładnością obwodów są to parametry składowe, a~nie błąd wyniku obwodu. Opracowanie własne.}
\label{tab:lr:noise:time}
\footnotesize
\setlength{\tabcolsep}{3pt}
\newcommand{\lrc}[1]{\raggedright\hspace{0pt}#1}
\begin{tabular}{@{}p{1.4cm}p{1.1cm}p{6.6cm}p{6.2cm}@{}}
\toprule
\lrc{\textbf{Skala}} & \lrc{\textbf{Praca}} & \lrc{\textbf{Urządzenie i~wielkość mierzona}} & \lrc{\textbf{Obserwacja}} \tabularnewline
\midrule
\lrc{minuty} & \lrc{\cite{hirasaki_detection_2023}} & \lrc{ibm\_kawasaki (27 kubitów): wynik pomiaru kubitu, seria 1000~s} & \lrc{skoki częstości błędów trwające po kilka minut} \tabularnewline
\lrc{godziny} & \lrc{\cite{burnett_decoherence_2019}} & \lrc{dwa transmony w~laboratorium: $\Tone$, serie 65 i~128~h} & \lrc{zmiany ponad dwukrotne; przełączenia co 2--12,5~h} \tabularnewline
\lrc{godziny} & \lrc{\cite{etxezarreta_multiqubit_2023}} & \lrc{pięć 5-kubitowych procesorów IBM: $\Tone$ co 2--4~min w~oknie kalibracji, serie 4--27~h} & \lrc{współczynnik zmienności 5--47\%} \tabularnewline
\lrc{godziny} & \lrc{\cite{kim_error_2025}} & \lrc{sześć transmonów: $\Tone$ (60~h) i~parametry uczonego modelu szumu (ok. 50~h)} & \lrc{$\Tone$ zmienne średnio o~ponad 300\%; skokowa zmiana parametrów modelu} \tabularnewline
\lrc{doba, dni} & \lrc{\cite{yeter-aydeniz_measuring_2023}} & \lrc{ibmq\_boeblingen: niewierność CNOT trzech par rano, wieczorem (po kalibracji bramek) i~po pięciu dniach} & \lrc{od $-13\%$ do $+30\%$ w~ciągu doby; wzrost 1,3--1,7 raza w~pięć dni (obliczenia własne)} \tabularnewline
\lrc{doba} & \lrc{\cite{wilson_just-in-time_2020}} & \lrc{procesory IBM: odczyt co godzinę; dokładność obwodów przy danych kalibracyjnych sprzed 28--49~h} & \lrc{odczyt niestabilny między kalibracjami; dokładność zmienna o~29--37\%} \tabularnewline
\lrc{miesiące} & \lrc{\cite{dasgupta_impact_2024}} & \lrc{IBM washington: odległość Hellingera rozkładów parametrów szumu, dane kalibracyjne z~16 miesięcy} & \lrc{0,41--0,92; stabilny wynik obwodu wymaga 0,022; największy wkład: pomiar} \tabularnewline
\bottomrule
\end{tabular}
\end{table}
% ŹRÓDŁA tabeli: hirasaki_detection_2023, streszczenie wersji opublikowanej (Crossref, DOI 10.1063/5.0166739) i opis eksperymentu (27 kubitów, seria 1000 s), arXiv:2307.04337; burnett_decoherence_2019, rys. 1 i 2 (ok. 65 h i ok. 128 h), ,,more than a factor 2'', czas przebywania 2--12,5 h, arXiv:1901.04417; etxezarreta_multiqubit_2023, tabele II i III, dodatek C, arXiv:2207.06838; kim_error_2025, sekcje II i III, arXiv:2407.02467v3; yeter-aydeniz_measuring_2023, tabela C3 (27 i 30 stycznia 2021 r.: zmiany +6,1; -7,1; +29,6 oraz -13,1; +4,6; -6,8 procent) i jej podpis, tabela C2 (24 i 29 stycznia: ilorazy 1,31; 1,61; 1,70), arXiv:2201.02899, obliczenia własne (PHD-TS-34); wilson_just-in-time_2020, sekcje II-A (rys. 1, IBM Poughkeepsie, 2019 r.) i V-B (rys. 6), arXiv:2005.12820; dasgupta_impact_2024, streszczenie, sekcje 3.3.1, 3.3.3 i 4, arXiv:2307.06833 (wynik z symulacji zasilanej danymi kalibracyjnymi; obwód Bernsteina--Vaziraniego na 5 kubitach).

Metody oparte na kalibracji reagują na tę zmienność ponawianiem kalibracji co okres kalibracji urządzenia albo co kilka godzin~\cite{maciejewski_mitigation_2020,van_den_berg_probabilistic_2023} (podrozdział~\ref{sec:lr:qem}); skokowe zmiany parametrów uczonego modelu szumu prowadziły według pracy~\cite{kim_error_2025} z~2025~r. do błędnych oszacowań po mitygacji.
% ŹRÓDŁO: jak w podrozdziale 3.3 (maciejewski_mitigation_2020, Remark 4, arXiv:1907.08518; van_den_berg_probabilistic_2023, materiał uzupełniający, część ,,Coherence'', arXiv:2201.09866); kim_error_2025, streszczenie i sekcja III, arXiv:2407.02467v3.
Sama częstość kalibracji procesorów IBM nie była stała: prace z~lat 2019--2023 podają dwie kalibracje na dobę~\cite{murali_noise-adaptive_2019}, okres 28--49~godzin bez kalibracji~\cite{wilson_just-in-time_2020}, kilka kalibracji na dobę dla urządzeń santiago i~bogota wobec około jednej dla quito, belem i~lima~\cite{etxezarreta_multiqubit_2023} oraz pełną kalibrację rano i~częściową wieczorem~\cite{yeter-aydeniz_measuring_2023}.
% ŹRÓDŁO: murali_noise-adaptive_2019, sekcja 2, arXiv:1901.11054; wilson_just-in-time_2020, sekcja V-B; etxezarreta_multiqubit_2023, dodatek C; yeter-aydeniz_measuring_2023, podpis tabeli C3.

\paragraph{Znaczenie dla założeń metody PDU.}
Wnioski dla PDU są pośrednie i~własne: prawie wszystkie przytoczone prace mierzą parametry składowe (czasy koherencji, błędy bramek i~odczytu, parametry modelu szumu), a~nie systematyczny błąd wyniku obwodu, którego dotyczą założenie~(\ref{itm:pdu:Z4}) i~warunek~\eqref{eq:pdu:stability}.
Z~tym zastrzeżeniem literatura nie potwierdza założenia~(\ref{itm:pdu:Z4}) w~postaci mocnej, czyli stałości błędu do następnej kalibracji urządzenia: przeciw niej świadczą zmiany skokowe w~skali minut, zmienność $\Tone$ w~oknie kalibracji, zmiany błędu bramek w~ciągu doby i~niestabilność odczytu.
Szum zmienny w~trakcie zbierania powtórzeń, których wymagają założenia~(\ref{itm:pdu:Z1}) i~(\ref{itm:pdu:Z3}), obciąża estymaty~\cite{kim_error_2025,proctor_detecting_2020}.
Literatura jest zgodna z~postacią słabszą: korekta wyznaczona eksperymentalnie ma czas ważności rzędu godzin, nie przenosi się między dniami ani miesiącami i~jest przypisana do konkretnych kubitów, bo fluktuacje są lokalne~\cite{burnett_decoherence_2019,etxezarreta_multiqubit_2023}, co odpowiada założeniu~(\ref{itm:pdu:Z5}).
Żadna ze sprawdzonych prac nie mierzy jednak przełożenia fluktuacji parametrów na dryf błędu wyniku, więc spełnienia warunku~\eqref{eq:pdu:stability} nie da się ocenić na podstawie literatury.
Zmiany epizodyczne może ponadto pominąć linia trendu błędu średniokwadratowego (podrozdział~\ref{subsec:pdu:results:stability}), a~reguła jednej kalibracji PDU na dobę (podrozdział~\ref{sec:pdu:procedure}) zakłada, że kalibracja urządzenia jest jednym, znanym użytkownikowi zdarzeniem, czego przytoczone obserwacje nie potwierdzają.
Zależność błędu od struktury programu i~historii operacji~\cite{proctor_measuring_2022,white_demonstration_2020} uzasadnia traktowanie obwodu jako całości (założenie~\ref{itm:pdu:Z2}), ale podważa przenoszenie korekty na obwody o~innej strukturze (założenie~\ref{itm:pdu:Z8}).

\paragraph{Szum a~symetrie: zaburzenia generatorów bramek.}\label{par:lr:noise:generators}
Z~symetriami wiąże szum praca z~2026~r.~\cite{raubitzek_quantum_noise}, której pierwszym autorem jest S.~Raubitzek, a~autor rozprawy jednym ze współautorów (implementował część kodu i~uruchamiał eksperymenty); nawiązuje ona do prac o~symetriach w~korekcji błędów~\cite{Faist_2020,liu2023quantumerrorcorrectionmeets,rajput2022quantumerrorcorrectiongauge,PhysRevA.97.032323}.
% ŹRÓDŁO udziału autora: ustalenie autora z 2026-11-26; oświadczenie o wkładzie autorów w wersji opublikowanej pracy, przed sekcją Data availability (Kr.W.: main manuscript, numerical simulation, experiments, and review); wersja zgłoszeniowa artykuły/quantum_noise_elsevier_submission/main.tex, wiersz 970.
% ŹRÓDŁO: rok i kolejność autorów: wpis raubitzek_quantum_noise w bibliography/own_publications.bib (APS Open Science 1, 000109, 2026, DOI 10.1103/hg15-cjnx; metadane potwierdzone w Crossref według recenzji eksperymenty/recenzje/rozdzial-9.md) i nagłówek wersji opublikowanej (published.txt, wiersze 1--14); nawiązanie do prac o symetriach: wersja opublikowana, sekcja VI.B (pozycje [44]--[47] wykazu literatury pracy).
% ŹRÓDŁO: Faist_2020, streszczenie, arXiv:1902.07714 (ciągłe symetrie ograniczają dokładność korekcji błędów); liu2023quantumerrorcorrectionmeets i rajput2022quantumerrorcorrectiongauge, streszczenia wersji opublikowanych (Crossref); PhysRevA.97.032323, streszczenie, arXiv:1709.05302 (prawa zachowania układu w wykrywaniu i korygowaniu błędów).
Sprawdza ona, czy małe addytywne zaburzenia generatorów bramek (podrozdział~\ref{subsec:tb:lie}), nazywane przez autorów łamiącymi symetrię, odtwarzają zrekonstruowaną na rzeczywistym urządzeniu macierz transferu Pauliego bramki (podrozdział~\ref{subsec:tb:noise:characterisation}).
% ŹRÓDŁO: raubitzek_quantum_noise, wersja opublikowana (eksperymenty/_build/PHD-TS-38/work/pub/published.txt), sekcja IV.A (każde zaburzenie: macierz zespolona 2x2 o ośmiu parametrach rzeczywistych, ograniczenie Delta_max = 0.05) i równanie (B42) (8|G| = 24); wersja zgłoszeniowa (artykuły/quantum_noise_elsevier_submission/main.tex), wiersze 418--432 i 483.
% ŹRÓDŁO: raubitzek_quantum_noise, wersja opublikowana, sekcja VI.A (,,They break exact unitarity, unitality, and trace preservation at the generator level'') i sekcja VII (,,symmetry-violating generator perturbations, especially non-Hermitian ones''); wersja zgłoszeniowa, wiersze 888 i 940.
Zaburzenia dopasowano algorytmem genetycznym do średniej macierzy z~88 niezależnych wykonań tomografii procesu bramki $\mathrm{X}$ na jednym kubicie procesora nadprzewodzącego ibm\_brisbane, przy czym zgodność każdego elementu oceniano względem rozrzutu pojedynczych przebiegów, a~nie błędu standardowego średniej.
% ŹRÓDŁO: raubitzek_quantum_noise, wersja opublikowana, sekcje III.A i IV (12 obwodów: cztery stany wejściowe i trzy bazy pomiarowe; zaburzenia przechodzą przez całe obwody tomografii, łącznie z przygotowaniem stanów i obrotami baz), sekcja IV (równanie (32) i zdanie po nim: średnia i odchylenie standardowe z próby po przebiegach), sekcja IV.A (równanie (35)), sekcja V.C (równanie (59): N_eval = 1000) i sekcja V.D (równania (68)--(71): moduł różnicy średnich odniesiony do pierwiastka z sumy kwadratów odchyleń standardowych z serii przebiegów na urządzeniu i z symulowanych rekonstrukcji); wersja zgłoszeniowa, wiersze 375--391, 498--503, 643 i 670--683.
% ŹRÓDŁO: raubitzek_quantum_noise, wersja opublikowana, sekcja IV (ibm_brisbane, Eagle r3; ,,For a chosen physical qubit''; ,,a fixed number of shots per circuit'') i równanie (42) (n_runs = 88); wersja zgłoszeniowa, wiersze 473 i 563.
Według pracy zbieranie danych trwało ponad pół roku, natomiast katalogi przebiegów w~jej publicznym repozytorium danych noszą daty dziewięciu dni między 3~lipca a~15~października 2025~r.
% ŹRÓDŁO: praca: raubitzek_quantum_noise, wersja opublikowana, sekcja VII (,,collecting the 88 hardware runs ... required more than half a year until IBM Brisbane was no longer available''); wersja zgłoszeniowa, wiersz 942. Repozytorium danych pracy (github.com/Raubkatz/QuantumSymmetryViolations, rewizja 5c4bd5dce7): wykaz plików zapisany w eksperymenty/_build/PHD-TS-38/work/repo/tree.json zawiera 88 katalogów przebiegów o nazwach z datami 20250703 (6 przebiegów), 20250716 (10), 20250807 (12), 20250808 (6), 20250809 (4), 20250915 (20), 20251007 (15), 20251013 (13) i 20251015 (2); policzone ponownie w zadaniu PHD-TS-94, zgodnie z recenzją (eksperymenty/recenzje/rozdzial-9.md, ustalenie I1).
Autorzy podają, że model z~dopasowanymi zaburzeniami odtwarza w~granicach tego kryterium większość pozadiagonalnych elementów średniej macierzy z~urządzenia, a~rozbieżności skupiają się na przekątnej bloku unitalnego.
% ŹRÓDŁO: raubitzek_quantum_noise, wersja opublikowana, sekcja VI, akapit trzeci (zgodność w granicach jednego łącznego odchylenia standardowego dla większości elementów; rozbieżności głównie w elementach (Y,Y) i (Z,Z)) i sekcja VII; zliczeń elementów zgodnych wersja opublikowana nie zawiera, a tekst ich nie podaje (ustalenie B1 recenzji).
Na tej przekątnej średnia z~urządzenia wynosi $0{,}904$, $-0{,}910$ i~$-0{,}908$, a~model daje $0{,}914$, $-0{,}966$ i~$-1{,}000$; różnice w~dwóch ostatnich elementach autorzy wiążą z~procesami dyssypacyjnymi, błędami przygotowania stanu i~pomiaru oraz dryfem, których model nie zawiera.
% ŹRÓDŁO: raubitzek_quantum_noise, wersja opublikowana, równania (45) (0,90412; -0,90994; -0,90794), (47) i (61) (0,91424; -0,96596; -0,99999); wersja zgłoszeniowa, wiersze 569--572 i 650--653; interpretacja autorów: sekcja VI, akapit trzeci, sekcja VII (,,missing dissipative structure or SPAM-related drift'') i dodatek B, punkt 10. Elementy przekątnej nie zależą od konwencji zapisu macierzy, której dotyczy ustalenie B1 recenzji.

Model autorzy określają jako ograniczony i~fenomenologiczny: ma 24 parametry rzeczywiste przy 16 elementach dopasowywanej macierzy, a~dopasowane zaburzenia są jedną z~możliwych parametryzacji obserwowanych odstępstw, a~nie wykazaną przyczyną fizyczną.
% ŹRÓDŁO: raubitzek_quantum_noise, wersja opublikowana, sekcja I (,,constrained phenomenological reconstruction model''), dodatek B, punkt 6 (równanie (B42) i zdanie po nim: 24 parametry rzeczywiste, 16 elementów macierzy), punkty 8 i 10 oraz sekcja VII (,,one possible effective parametrization ... rather than ... a unique microscopic physical noise mechanism''); w wersji zgłoszeniowej tych zastrzeżeń nie ma.
Badanie obejmuje jeden kubit, jedną bramkę i~jedno urządzenie, na którym pomiarów nie da się już powtórzyć, a~zrekonstruowane macierze mogą, jak zaznaczają autorzy, zawierać także wkład błędów przygotowania stanu i~pomiaru, dryfu kalibracji i~skończonej liczby powtórzeń.
% ŹRÓDŁO: raubitzek_quantum_noise, wersja opublikowana, streszczenie i sekcja VI (ograniczenie do bramek jednokubitowych), sekcje IV i VII (jeden kubit, bramka X, ibm_brisbane; urządzenie niedostępne).
% ŹRÓDŁO: raubitzek_quantum_noise, wersja opublikowana, sekcja III.A (procedura nie jest odporna na błędy SPAM), sekcja VI (,,may also contain contributions from SPAM effects, finite-shot fluctuations, calibration drift'') i dodatek B, punkt 8.
Zaburzenia dopasowano raz, do średniej ze wszystkich przebiegów, a~praca nie rozdziela źródeł rozrzutu między przebiegami; jej streszczenie mówi o~stabilnych zaburzeniach w~dopasowanych modelach, lecz ich stabilności w~czasie w~pracy nie mierzono.
% ŹRÓDŁO: raubitzek_quantum_noise, wersja opublikowana, streszczenie (,,the appearance of stable symmetry-violating perturbations in the fitted models''), sekcja I (uśrednianie po przebiegach dla stłumienia szumu statystycznego i dryfu), sekcje IV.A i V.B (jedno dopasowanie do średniej, ,,best-found perturbations''); wersja zgłoszeniowa, wiersz 70; por. recenzja, ustalenie I2.
Naszkicowanego w~pracy schematu wykorzystania takiego opisu w~mitygacji i~korekcji błędów nie zrealizowano i~nie skorygowano w~niej wyniku żadnego obliczenia; autorzy zastrzegają, że model może obejmować tylko część szumu i~że schemat wymaga zaburzeń stabilnych, powtarzalnych i~mierzalnych w~skali czasu kodowania i~korekcji, co odpowiada założeniu~(\ref{itm:pdu:Z4}) metody PDU (podrozdział~\ref{sec:pdu:discussion}).
% ŹRÓDŁO: raubitzek_quantum_noise, wersja opublikowana, sekcja VI.B (trzy kroki: scharakteryzować; uzupełnić etap kodowania albo odtwarzania o operacje usuwające rozpoznane odstępstwa; monitorować ,,the first row and affine translation of the PTM'' jako wskaźniki pozostałego łamania symetrii); wersja zgłoszeniowa, wiersze 911--918. Sformułowanie nie zależy od konwencji zapisu (ustalenie B1 recenzji).
% ŹRÓDŁO: raubitzek_quantum_noise, wersja opublikowana, sekcja VI.B, zastrzeżenia (1) i (2) oraz zdanie końcowe o przyszłych badaniach; wersja zgłoszeniowa, wiersze 911--918.
\agenttask{
Pytania do autora (PHD-TS-94; punkty 2--5 przeniesione z dawnej sekcji 9.2, zadanie PHD-TS-38).
(1) Ustalenie B1 recenzji dawnego rozdziału 9 \agentref{eksperymenty/recenzje/rozdzial-9.md}: według recenzenta średnia macierz transferu Pauliego z urządzenia (równanie 45 pracy) jest zapisana w konwencji transponowanej względem definicji pracy i względem symulatora, więc to, co praca nazywa wierszem afinicznym urządzenia, może być wektorem przesunięcia, a liczba elementów spełniających kryterium zgodności może po ujednoliceniu zapisu być kilkakrotnie mniejsza od podanej w wersji zgłoszeniowej pracy (liczby w recenzji). Czy autor i współautorzy to potwierdzają? Sprawdzenie według recenzji: porównać qiskit.quantum_info.PTM(choi).data z wynikiem wzoru ze skryptu pomiarowego dla dowolnego zapisanego przebiegu albo uruchomić skrypt zliczenia_a_macierz.py z katalogu eksperymenty/recenzje/_robocze/PHD-TS-66. Do czasu odpowiedzi akapit nie podaje zliczeń elementów zgodnych ani wartości wskaźnika zgodności, nie przytacza elementów macierzy spoza przekątnej bloku unitalnego i nie rozstrzyga, która część macierzy jest wierszem afinicznym, a która wektorem przesunięcia; po odpowiedzi uzupełnić akapit i sprawdzić zdanie o wierszu afinicznym w podrozdziale 2.2 (theoretical_background_02_noise.tex).
(2) Wersja opublikowana pracy (PDF pobrany z harvest.aps.org, 49 stron) nie zawiera tabel wskaźnika zgodności, macierzy zespołów losowych ani zliczeń elementów zgodnych, które są w wersji zgłoszeniowej; materiałów dodatkowych wydawcy nie sprawdzano. Czy tabele usunięto w recenzji?
(3) Okres zbierania danych: praca podaje ponad pół roku, a katalogi 88 przebiegów w publicznym repozytorium pracy mają daty dziewięciu dni od 2025-07-03 do 2025-10-15 (około 3,5 miesiąca). Tekst podaje obie informacje ze wskazaniem źródła; która jest właściwa?
(4) Praca nie podaje liczby powtórzeń obwodu na urządzeniu ani numeru kubitu; czy autor chce je podać z danych (według recenzji 10 000 powtórzeń w 10 sprawdzonych plikach zliczeń i w konfiguracji skryptu)?
(5) Recenzja (ustalenie I1) proponuje mocniejsze sformułowanie o zmienności oparte na ponownym odczytaniu publicznych danych pracy (udział wariancji między dniami pomiarowymi); nie wpisano go, bo w części dotyczy elementów objętych ustaleniem B1 i jest obliczeniem własnym spoza pracy. Czy je dodać po rozstrzygnięciu punktu (1)?
(6) Akceptacja sformułowań, które nie są twierdzeniami autorów pracy: że kryterium zgodności odnosi różnicę do rozrzutu pojedynczych przebiegów, a nie do błędu standardowego średniej; że streszczenie mówi o stabilnych zaburzeniach, a praca nie zawiera pomiaru stabilności; że warunek stabilności schematu odpowiada założeniu Z4.
}

